\documentclass[11pt]{article}

\usepackage[margin=0.85in]{geometry}
\usepackage{amsmath}
\usepackage{array}
\usepackage{booktabs}
\usepackage{hyperref}
\usepackage{xcolor}
\usepackage{enumitem}

\hypersetup{
    colorlinks=true,
    linkcolor=blue,
    urlcolor=blue
}

\setlength{\parindent}{0pt}
\setlength{\parskip}{0.65em}
\renewcommand{\arraystretch}{1.18}

\title{Lissajous Ground-Scan Implementation}
\author{SimGroundLissajous}
\date{}

\begin{document}
\maketitle

\section*{Purpose}

The \texttt{SimGroundLissajous} operator simulates a smooth ground-based
Lissajous boresight scan for a TOAST observation.  It follows the same broad
workflow as a ground-scan operator: it reads a ground schedule, creates
observations, generates timestamps and boresight pointing, converts the
mount-frame pointing to sky coordinates, and records useful observation
metadata.

The important design choice is that the scan is generated natively in the
mount coordinate system: azimuth and elevation.  The equatorial pointing is
then derived afterward by converting the generated Az/El boresight at each
timestamp.

\section*{Coordinate Picture}

The scan is first described in horizontal coordinates:

\begin{itemize}[leftmargin=*]
    \item Azimuth, \(\mathrm{Az}\), is the commanded horizontal mount angle.
    \item Elevation, \(\mathrm{El}\), is the commanded elevation mount angle.
    \item RA/Dec are produced later by converting the time-dependent Az/El
    boresight through the observing site and timestamp.
\end{itemize}

So the operator is not scanning directly in equatorial coordinates.  It scans
in Az/El and then asks: ``where on the sky was this mount pointing at this
time?''

\section*{Schedule Rows and Observations}

Each schedule row supplies a time interval and a nominal mount pointing range:

\begin{center}
\begin{tabular}{>{\raggedright\arraybackslash}p{0.26\linewidth}
                >{\raggedright\arraybackslash}p{0.60\linewidth}}
\toprule
Schedule quantity & Role in the Lissajous scan \\
\midrule
Start / stop time & Defines the time span of one generated scan. \\
\(\mathrm{Az}_{\min}\), \(\mathrm{Az}_{\max}\) & Define the default azimuth center and, optionally, the default azimuth amplitude. \\
\(\mathrm{El}\) & Defines the default elevation center. \\
Rotation / boresight angle & Applied as a boresight rotation after the Az/El pointing is generated. \\
\bottomrule
\end{tabular}
\end{center}

For each schedule row, the operator creates one Lissajous scan over that row's
time interval.  With the default zero drift rates, the scan center is fixed in
mount coordinates during that observation.  The sky moves through that fixed
mount pattern as Earth rotates.

At the next schedule row, the operator recomputes the center from the new
scheduled \(\mathrm{Az}_{\min}\), \(\mathrm{Az}_{\max}\), and \(\mathrm{El}\),
then starts a new Lissajous pattern.

\section*{Core Scan Equations}

For each observation, the operator defines a local elapsed time

\[
    t = \frac{i}{f_{\mathrm{samp}}},
\]

where \(i\) is the sample index and \(f_{\mathrm{samp}}\) is the detector sample
rate.

The azimuth and elevation centers are

\[
    \mathrm{Az}_{0} =
    \frac{\mathrm{Az}_{\min} + \mathrm{Az}_{\max}}{2},
    \qquad
    \mathrm{El}_{0} = \mathrm{El}_{\mathrm{schedule}}.
\]

The Lissajous trajectory is then

\[
    \mathrm{Az}(t) =
    \mathrm{Az}_{0}
    + v_{\mathrm{Az,drift}} t
    + A_{\mathrm{Az}}
      \sin \left(2\pi f_{\mathrm{Az}} t + \phi_{\mathrm{Az}}\right),
\]

\[
    \mathrm{El}(t) =
    \mathrm{El}_{0}
    + v_{\mathrm{El,drift}} t
    + A_{\mathrm{El}}
      \sin \left(2\pi f_{\mathrm{El}} t + \phi_{\mathrm{El}}\right).
\]

Here \(A\) is amplitude, \(f\) is oscillation frequency, \(\phi\) is phase, and
\(v_{\mathrm{drift}}\) is an optional linear drift of the scan center.

\section*{What Zero Drift Means}

By default,

\[
    v_{\mathrm{Az,drift}} = 0,
    \qquad
    v_{\mathrm{El,drift}} = 0.
\]

This means the Lissajous pattern is fixed around the scheduled Az/El center
for the duration of one observation.  The telescope mount repeats the same
mount-frame pattern, while the sky drifts through it.

In plain language:

\begin{quote}
Within one observation, hold the scheduled mount center and draw a Lissajous
around it.  At the next schedule row, use the newly scheduled center and draw a
new Lissajous pattern.
\end{quote}

This is discrete re-centering between observations, not continuous source
tracking inside a single observation.

\section*{What Nonzero Drift Does}

The drift terms move the center of the Lissajous pattern linearly during an
observation.  For example, a positive elevation drift raises the entire pattern
in elevation as time passes.

These drift rates can be used to approximate tracking, to deliberately sweep a
larger region, or to move the pattern from one scheduled center toward another.
They do not automatically compute the source tracking rate, and they do not
guarantee that the same RA/Dec path is retraced.

A nonzero drift usually makes the mount path an ``open'' Lissajous sweep rather
than a closed repeating figure centered at one fixed Az/El position.

\section*{Default Amplitudes and Centers}

The azimuth amplitude has a schedule-aware default:

\[
    A_{\mathrm{Az}} =
    \frac{\left|\mathrm{Az}_{\max} - \mathrm{Az}_{\min}\right|}{2}
    \quad \mathrm{if}\quad
    \texttt{lissajous\_az\_amplitude = None}.
\]

If \texttt{lissajous\_az\_amplitude} is explicitly set, that value is used
instead.

The elevation amplitude defaults to the configured trait value.  It is not
derived from the schedule because the schedule provides a single elevation
center rather than an elevation min/max range.

\section*{Mount Motion Limit Checks}

The operator checks the requested Lissajous motion against user supplied mount
motion limits.  These limits are mechanical mount-axis limits, not on-sky scan
speed limits.

The peak azimuth and elevation mount rates are computed analytically:

\[
    \left|\dot{\mathrm{Az}}\right|_{\max}
    =
    \left|v_{\mathrm{Az,drift}}\right|
    + 2\pi f_{\mathrm{Az}} \left|A_{\mathrm{Az}}\right|,
\]

\[
    \left|\dot{\mathrm{El}}\right|_{\max}
    =
    \left|v_{\mathrm{El,drift}}\right|
    + 2\pi f_{\mathrm{El}} \left|A_{\mathrm{El}}\right|.
\]

The peak accelerations are

\[
    \left|\ddot{\mathrm{Az}}\right|_{\max}
    =
    \left(2\pi f_{\mathrm{Az}}\right)^2
    \left|A_{\mathrm{Az}}\right|,
\]

\[
    \left|\ddot{\mathrm{El}}\right|_{\max}
    =
    \left(2\pi f_{\mathrm{El}}\right)^2
    \left|A_{\mathrm{El}}\right|.
\]

These peak values are compared to
\texttt{max\_mount\_rate\_az}, \texttt{max\_mount\_rate\_el},
\texttt{max\_mount\_accel\_az}, and \texttt{max\_mount\_accel\_el}.

Because the scan is generated in Az/El, this is a mount-coordinate check:

\[
    \dot{\mathrm{Az}} \quad \mathrm{and} \quad \dot{\mathrm{El}}.
\]

It is not an on-sky angular speed check.  A small-angle on-sky speed would look
approximately like

\[
    v_{\mathrm{sky}}
    \simeq
    \sqrt{
        \left(\cos(\mathrm{El})\,\dot{\mathrm{Az}}\right)^2
        +
        \dot{\mathrm{El}}^2
    }.
\]

That \(\cos(\mathrm{El})\) projection is intentionally not applied in the mount
limit check.

\section*{Main User-Facing Traits}

\begin{center}
\begin{tabular}{>{\ttfamily}p{0.35\linewidth}
                >{\raggedright\arraybackslash}p{0.52\linewidth}}
\toprule
Trait & Meaning \\
\midrule
lissajous\_az\_amplitude & Azimuth oscillation amplitude.  If \texttt{None}, use half the scheduled azimuth throw. \\
lissajous\_el\_amplitude & Elevation oscillation amplitude. \\
lissajous\_az\_rate & Azimuth oscillation frequency. \\
lissajous\_el\_rate & Elevation oscillation frequency. \\
lissajous\_az\_phase & Azimuth sinusoid phase offset. \\
lissajous\_el\_phase & Elevation sinusoid phase offset. \\
lissajous\_az\_drift\_rate & Linear drift rate of the azimuth scan center. \\
lissajous\_el\_drift\_rate & Linear drift rate of the elevation scan center. \\
check\_mount\_motion\_limits & Enable or disable mount rate and acceleration checks. \\
max\_mount\_rate\_az & Maximum allowed azimuth mount rate. \\
max\_mount\_rate\_el & Maximum allowed elevation mount rate. \\
max\_mount\_accel\_az & Maximum allowed azimuth mount acceleration. \\
max\_mount\_accel\_el & Maximum allowed elevation mount acceleration. \\
\bottomrule
\end{tabular}
\end{center}

\section*{Operational Traits}

\begin{center}
\begin{tabular}{>{\ttfamily}p{0.35\linewidth}
                >{\raggedright\arraybackslash}p{0.52\linewidth}}
\toprule
Trait & Meaning \\
\midrule
telescope & Telescope object with focalplane and site information. \\
telescope\_file & Optional instrument file used when a telescope object is not supplied. \\
schedule & GroundSchedule object used to define observation rows. \\
schedule\_file & Optional schedule file used when a schedule object is not supplied. \\
sort\_schedule\_file & Sort a loaded schedule by mean boresight RA. \\
weather & Built-in weather site name or weather HDF5 file. \\
realization & Weather realization index. \\
max\_pwv & Optional maximum PWV for simulated weather. \\
median\_weather & Use median weather parameters instead of sampling distributions. \\
use\_ephem & Use PyEphem for Az/El to RA/Dec conversion. \\
use\_qpoint & Use qpoint for Az/El to RA/Dec conversion. \\
distribute\_time & Distribute observations along time rather than detector axis. \\
sample\_set\_size & Samples per sample-set when distributing in time. \\
detset\_key & Focalplane column used to group detectors into detector sets. \\
\bottomrule
\end{tabular}
\end{center}

\section*{Produced Data}

For each observation, the operator writes shared data columns for:

\begin{itemize}[leftmargin=*]
    \item timestamps,
    \item azimuth,
    \item elevation,
    \item boresight quaternion in Az/El,
    \item boresight quaternion in RA/Dec,
    \item observatory position,
    \item observatory velocity,
    \item optional shared flags.
\end{itemize}

It also records scan metadata such as the realized minimum and maximum Az/El
values and marks the scan type as \texttt{lissajous}.

\section*{What This Implementation Does Not Do}

\begin{itemize}[leftmargin=*]
    \item It does not continuously solve for source tracking inside each
    observation.
    \item It does not automatically compute drift rates from a target's
    apparent Az/El motion.
    \item It does not check absolute mechanical pointing bounds such as
    \(\mathrm{Az}_{\min}^{\mathrm{mount}}\) or
    \(\mathrm{El}_{\min}^{\mathrm{mount}}\).
    \item It does not check total on-sky angular speed.
    \item It does not generate physical turnaround periods.  For compatibility
    with TOAST operators that expect the standard key, it creates an empty
    \texttt{turnaround} interval list.
\end{itemize}

\section*{Practical Mental Model}

The simplest way to understand the current implementation is:

\begin{quote}
The schedule tells the telescope where to place the center of each pass in
Az/El.  The Lissajous traits tell it how to wiggle around that center.  With
zero drift, the sky passes through that mount-fixed pattern.  At the next
schedule row, the pattern is re-centered and the process repeats.
\end{quote}

This makes the operator useful for simulating compact, smooth, repeated
mount-frame patterns while still producing the RA/Dec pointing needed by later
mapmaking and sky-projection steps.

\end{document}
